\Subhead{5.2.3}{Распространение радиоволн в ионосфере. Дисперсия, поглощение и рефракция}

Ионосфера содержит свободные электроны; её свойства зависят от высоты, освещения и состояния плазмы. В простой немагнитной модели со столкновениями частоты $\nu_c$
\begin{equation}
 \varepsilon_r(\omega,h)=1-\frac{\omega_p^2(h)}{\omega(\omega+\ii\nu_c)},
 \qquad \omega_p^2=\frac{n_ee^2}{m_e\varepsilon_0}.
\end{equation}
При малых столкновениях $n^2\simeq1-\omega_p^2/\omega^2$. Фазовая скорость больше $c$, групповая меньше $c$; передача новой информации остаётся причинной. Столкновения создают положительную мнимую часть $n$ и поглощение, особенно существенное там, где плотность нейтрального газа велика. Геомагнитное поле расщепляет нормальные волны и поляризации; тензорная модель — в вопросе~\Qref{2.3}.

В плоско-слоистой немагнитной ионосфере сохраняется $p=n\sin\theta$. Вертикальный луч поворачивает около уровня $\omega=\omega_p(h)$; максимум плазменной частоты слоя задаёт критическую частоту вертикального зондирования. Для наклонного падения и $n\simeq1$ у земли условие поворота даёт секансный закон $f\simeq f_p(h_t)/\cos\theta_0$. Он объясняет, почему наклонная трасса использует частоты выше вертикальной критической, но не учитывает кривизну Земли, магнитное расщепление и реальные профили.

Луч поворачивает постепенно вследствие градиента показателя; описание «отражение от слоя» — приближённый язык. На земле возможна зона отсутствия возвращающегося ионосферного луча; конечное поглощение и неоднородности меняют дальность и замирания.
